\chapter{Estimarea matricei de precizie prin graphical lasso}

Dupa eliminarea dependentei temporale prin modelul VAR, atentia se muta asupra
structurii de dependenta contemporana dintre reziduuri. Obiectivul este estimarea
matricei de precizie $\Theta = \Sigma^{-1}$, ale carei elemente nenule definesc
reteaua financiara.

\section{Limitele estimarii directe}

O prima idee ar fi sa estimam $\Theta$ prin inversarea matricei de covarianta
empirica $\hat{\Sigma}$. Aceasta abordare are insa limitari serioase:

\begin{enumerate}
    \item \textit{Instabilitate numerica}: cand $n$ este de aceeasi ordine cu
    $T$, $\hat{\Sigma}$ poate fi prost conditionata, iar $\hat{\Sigma}^{-1}$
    amplifica erorile de estimare.
    \item \textit{Lipsa de sparsitate}: inversa empirica este in general densa,
    producand un graf complet greu de interpretat.
    \item \textit{Supraajustare}: o matrice de precizie densa include inevitabil
    relatii slabe sau accidentale care nu reflecta structura reala a datelor.
\end{enumerate}

In contextul acestei lucrari, cu $n = 11$ sectoare si $T \approx 1900$ observatii,
raportul $n/T \approx 0.006$ este favorabil, dar sparsitatea este totusi
dorita pentru interpretabilitate. Metoda graphical lasso rezolva ambele probleme.

\section{Problema de optimizare penalizata}

Metoda graphical lasso (GLasso), propusa de \cite{friedman2008}, estimeaza
matricea de precizie prin maximizarea log-verosimilitatii penalizate:

\begin{equation}
    \hat{\Theta} = \underset{\Theta \succ 0}{\arg\max}
    \Bigl[\log\det(\Theta) - \operatorname{tr}(\hat{S}\,\Theta)
    - \rho\,\|\Theta\|_1\Bigr],
    \label{eq:glasso}
\end{equation}

unde:
\begin{itemize}
    \item $\hat{S}$ este matricea de covarianta empirica a reziduurilor VAR,
    \item $\rho > 0$ este parametrul de regularizare,
    \item $\|\Theta\|_1 = \sum_{i \neq j} |\theta_{ij}|$ este norma $\ell_1$
    aplicata elementelor extra-diagonale (diagonala nu este penalizata).
\end{itemize}

\subsection{Interpretarea termenilor}

Termenul $\log\det(\Theta) - \operatorname{tr}(\hat{S}\,\Theta)$ reprezinta
log-verosimilitatea gaussiana (pana la o constanta) in functie de $\Theta$.
Maximizarea sa fara penalizare ar produce $\hat{\Theta} = \hat{S}^{-1}$,
adica inversa empirica. Termenul de penalizare $-\rho\|\Theta\|_1$ forteaza
elementele extra-diagonale mici spre zero, producand o matrice rara.

\subsection{Convexitate si existenta solutiei}

\begin{proposition}
Problema de optimizare \eqref{eq:glasso} este strict concava in $\Theta \succ 0$
si admite o solutie unica $\hat{\Theta}$.
\end{proposition}

\begin{proof}
Functia obiectiv este suma a doua componente: $\log\det(\Theta)$, strict concava
pe conul pozitiv definit, si $-\operatorname{tr}(\hat{S}\,\Theta) - \rho\|\Theta\|_1$,
concava (liniara plus concava). Prin urmare, suma este strict concava, iar
maximul pe multimea inchisa $\{\Theta : \Theta \succ 0\}$ este unic.
\end{proof}

\subsection{Algoritm de rezolvare}

Problema \eqref{eq:glasso} este rezolvata prin algoritmul de coordonate-ciclice
propus de \cite{friedman2008}. Ideea este de a optimiza alternant in fiecare
coloana/linie a lui $\Theta$, mentinand celelalte fixe. Fiecare subproblema
se reduce la un lasso standard, rezolvabila eficient. Convergenta algoritmului
la solutia globala este garantata de convexitatea problemei.

In R, metoda este implementata prin functia \texttt{glasso()} din pachetul
\texttt{glasso} \cite{friedman2019}, apelata in scriptul
\texttt{01\_var\_glasso\_network.R}.

\section{Selectia parametrului de regularizare}

Parametrul $\rho$ controleaza fundamental rezultatul:
\begin{itemize}
    \item $\rho \to 0$: $\hat{\Theta} \to \hat{S}^{-1}$, graf complet.
    \item $\rho \to \infty$: $\hat{\Theta} \to \operatorname{diag}(\hat{S})^{-1}$,
    graf fara nicio muchie.
\end{itemize}

Alegerea valorii optime $\hat{\rho}$ este realizata in scriptul
\texttt{03\_glasso\_selection.R} prin doua criterii complementare.

\subsection{Criteriul BIC pentru graphical models}

Criteriul BIC adaptat pentru modele grafice gaussiene \cite{yuan2007} este:
\begin{equation}
    \mathrm{BIC}(\rho) = -T\Bigl[\log\det(\hat{\Theta}_\rho)
    - \operatorname{tr}(\hat{S}\,\hat{\Theta}_\rho)\Bigr]
    + \ln(T) \cdot |\hat{E}_\rho|,
    \label{eq:bic_glasso}
\end{equation}
unde $|\hat{E}_\rho|$ este numarul de elemente extra-diagonale nenule ale lui
$\hat{\Theta}_\rho$ (numarul de muchii estimate). Valoarea optima este:
\[
    \hat{\rho} = \arg\min_{\rho > 0} \mathrm{BIC}(\rho).
\]

\subsection{Analiza pe grila de valori}

Selectia este realizata pe o grila logaritmica de valori
$\rho \in \{0.001, 0.002, \ldots, 0.5\}$. Pentru fiecare valoare, se estimeaza
$\hat{\Theta}_\rho$ si se calculeaza $\mathrm{BIC}(\rho)$ si densitatea retelei:
\[
    \mathrm{densitate}(\rho) = \frac{|\hat{E}_\rho|}{n(n-1)/2}.
\]

Compromisul dintre fidelitatea fata de date si parsimonia modelului este
esential: o retea excesiv de densa este greu de interpretat, iar una excesiv
de rara poate pierde relatii relevante.

\section{Proprietati statistice ale estimatorului GLasso}

\begin{proposition}[Consistenta, \cite{ravikumar2011}]
Sub conditii de regularitate asupra structurii reale a lui $\Theta^*$ si cu
$\rho \sim \sqrt{\ln n / T}$, estimatorul GLasso satisface:
\[
    \|\hat{\Theta}_\rho - \Theta^*\|_F = O_P\!\left(\sqrt{\frac{n \ln n}{T}}\right),
\]
si recupereaza structura corecta a grafului cu probabilitate tending la 1.
\end{proposition}

Aceasta proprietate garanteaza ca, pentru un numar suficient de mare de
observatii, structura retelei estimate converge la structura adevarata a
dependentelor conditionale.

\section{De la matricea de precizie la retea}

Odata estimata $\hat{\Theta}$, reteaua financiara este definita astfel:
\begin{itemize}
    \item \textbf{Noduri}: $V = \{1, \ldots, 11\}$, corespunzand celor 11 sectoare.
    \item \textbf{Muchii}: $\{i,j\} \in E \iff \hat{\theta}_{ij} \neq 0$ si $i \neq j$.
    \item \textbf{Ponderi}: $w_{ij} = |\hat{\theta}_{ij}|$, reprezentand intensitatea
    dependentei conditionale.
\end{itemize}

Prin constructie, aceasta retea nu reflecta corelatii brute, ci dependente
conditionale: o muchie intre sectoarele $i$ si $j$ indica o legatura directa
care persista dupa controlul pentru toate celelalte sectoare. Aceasta proprietate
este cea care face reteaua informativa din punct de vedere economic.

\section{Selectia parametrului de regularizare}

Un element central al metodei graphical lasso este alegerea parametrului de
regularizare $\rho$. Acesta controleaza compromisul dintre fidelitatea fata de
date si raritatea structurii estimate. Pentru valori mici ale lui $\rho$,
matricea de precizie estimata tinde sa fie mai densa, iar reteaua rezultata
devine mai greu de interpretat. Pentru valori mari ale lui $\rho$, multe muchii
sunt eliminate, iar reteaua devine mai rara.

In implementarea lucrarii, sensibilitatea fata de acest parametru este analizata
pe o grila de valori, iar pentru fiecare valoare sunt evaluate densitatea
retelei si un criteriu de tip BIC. Aceasta analiza are rolul de a arata cum se
modifica structura estimata atunci cand regularizarea devine mai puternica.

Din punct de vedere statistic, alegerea lui $\rho$ nu este doar o decizie
numerica, ci una interpretativa. O retea prea densa poate masca relatiile
esentiale printr-un exces de muchii slabe, in timp ce o retea prea rara poate
omite legaturi conditionale reale.

\section{Legatura cu implementarea empirica}

In analiza principala si in analiza pe ferestre temporale se foloseste aceeasi
valoare a penalizarii, pentru a pastra comparabilitatea intre rezultatele
obtinute. Astfel, modificarile observate in structura retelei de la o fereastra
la alta pot fi atribuite mai usor schimbarilor din date si nu unei modificari
a parametrului de regularizare.

Aceasta decizie metodologica este importanta in special in capitolul dedicat
ferestrelor temporale, unde interesul principal este compararea topologiei
retelei in timp, nu recalibrarea separata a modelului pentru fiecare subesantion.